% Môn: Vật lý
% Lớp: 11
% Chương: Chương IV. Dòng điện. Mạch điện
% Bài: L11C4  Bài 25. Năng lượng điện và công suất điện
% Số câu: 24
% App đọc file này trực tiếp — sửa rồi Commit trên GitHub, bấm Đồng bộ GitHub .tex

% L11C4  Bài 25. Năng lượng điện và công suất điện

% ===== Câu 1 =====
% ID: L11C4B25-01-TL
% Mức: NB
\dangbt{Chưa có dạng}
\begin{ex}
% ID: L11C4B25-01-TL
% Mức: NB
Hai điện trở R_1 và R_2 được mắc vào hiệu điện thế không đổi. Hỏi trong trường hợp mắc nối tiếp hai điện trở và mắc song song hai điện trở thì đoạn mạch nào tiêu thụ công suất lớn hơn? BÀ $I$ TẬ $P$ CUỐ $I$ CHƯƠNG IV $\nguon{ly11 kntt.pdf}$
\choiceTF

\loigiai{
Khi hai điện trở mắc nối tiếp, công suất tiêu thụ là: $\text{P}_{1} = \dfrac{U^{2}}{R_{1} + R_{2}}.$

Khi hai điện trở mắc song song, công suất tiêu thụ là: $\text{P}_{2} = \dfrac{U^{2}\left( {R_{1} + R_{2}} \right)}{R_{1}R_{2}}$

Từ hai biểu thức trên ta thấy, công suất khi hai điện trở mắc song song lớn hơn khi hai điện trở mắc nối tiếp vì điện trở tương đương của hai điện trở khi mắc song song nhỏ hơn khi mắc nối tiếp.

Ta có: $\dfrac{\text{P}_{2}}{\text{P}_{1}} = \dfrac{\left( {R_{1} + R_{2}} \right)^{2}}{R_{1}R_{2}}$.

Theo định lí Cauchy cho hai số dương 1 và 2 thì $R_{1} + R_{2} \geq 2\sqrt{R_{1}R_{2}}$ nên ta có:

$\left. \dfrac{\text{P}_{2}}{\text{P}_{1}} \geq \dfrac{4\left( \sqrt{R_{1}R_{2}} \right)^{2}}{R_{1}R_{2}}\Rightarrow\dfrac{\text{P}_{2}}{\text{P}_{1}} \geq 4. \right.$

Như vậy khi hai điện trở mắc song song thì công suất lớn hơn.
}
\end{ex}

% ===== Câu 2 =====
% ID: L11C4B25-01-TN
% Mức: NB
\begin{ex}
% ID: L11C4B25-01-TN
% Mức: NB
Đơn vị đo năng lượng điện tiêu thụ là $\nguon{ly11 kntt.pdf}$
\choice
{$\mathrm{kW}.$}
{kV.}
{k $\Omega$.}
{\True $\mathrm{kW}.h.$}
\loigiai{
Chọn D.

Đơn vị đúng kW.h.
}
\end{ex}

% ===== Câu 3 =====
% ID: L11C4B25-02-TL
% Mức: TH
\begin{ex}
% ID: L11C4B25-02-TL
% Mức: TH
Một bếp điện hoạt động liên tục trong 4 giờ ở hiệu điện thế $220\,\mathrm{V}$. Khi đó, số chỉ của công tơ điện tăng thêm $3\,\mathrm{s}$ ố. Công suất tiêu thụ của bếp điện và cường độ dòng điện chạy qua bếp trong thời gian trên là bao nhiêu? A. $P =$ 750 $\mathrm{kW}$ và I = 341\ $,$ \mathrm{A} $. B.$ P = 750\ $,$ \mathrm{W} $và$ I = 3 $,$ 41\,\mathrm{A} $. C.$ P = 750\ $,$ \mathrm{J} $và$ I = 3 $,$ 41\,\mathrm{A} $. D.$ P = 750\ $,$ \mathrm{W} $và$ I = 3 $.$ 14\,\mathrm{A} $.$ \nguon{ly11kntt.pdf}
\choiceTF

\loigiai{
Chọn B.

Công suất tiêu thụ $\text{P} = \dfrac{A}{t} = \dfrac{3kWh}{4h} = 0,$ 75 $\mathrm{kW}$

Cường độ dòng điện $I$ = $\dfrac{\text{P}}{U}$ = $\dfrac{750}{220}$ = 3, $41\,\mathrm{A}$
}
\end{ex}

% ===== Câu 4 =====
% ID: L11C4B25-02-TN
% Mức: NB
\begin{ex}
% ID: L11C4B25-02-TN
% Mức: NB
Cho dòng điện $I$ chạy qua hai điện trở R_1 và R_2 mắc nối tiếp. Mối liên hệ giữa nhiệt lượng toả ra trên mỗi điện trở và giá trị các điện trở là: $QRQR\nguon{ly11 kntt.pdf}$
\choice
{\True $1 = 1.$}
{$2 = 1.$}
{Q_1 $R_1 = Q_2R 2$.}
{Q_1 $Q_1 = R_1R 2 Q_2 R 2 Q_1 R 2$}
\loigiai{
Chọn A.

Ta có
}
\end{ex}

% ===== Câu 5 =====
% ID: L11C4B25-03-TL
% Mức: TH
\begin{ex}
% ID: L11C4B25-03-TL
% Mức: TH
Một trường học có 20 phòng học, tính trung bình mỗi phòng học sử dụng điện trong 10 giờ mỗi ngày với một công suất điện tiêu thụ $500\,\mathrm{W}$. a) Tính công suất điện tiêu thụ trung bình của trường học trên. b) Tính năng lượng điện tiêu thụ của trường học trên 30 ngày. c) Tính tiền điện của trường học trên phải trả trong 30 ngày với giá điện 2000 đ/ $\mathrm{kW}$.h. d) Nếu tại các phòng học của trường học trên, các bạn học $\sin$ h đều có ý thức tiết kiệm điện bằng cách tắt các thiết bị điện khi không sử dụng. Thời gian dùng các thiết bị điện ở mỗi phòng học chỉ còn 8 giờ mỗi ngày. Em hãy tính tiền điện mà trường học trên đã tiết kiệm được trong một năm học (9 tháng, mỗi tháng $30$ ngày). $\nguon{ly11 kntt.pdf}$
\choiceTF

\loigiai{
a) Công suất điện tiêu thụ trung bình của trường học:

$P= 500.20 =10000\,\mathrm{W}=10\,\mathrm{kW}$.

b) Năng lượng điện tiêu thụ của trường học trong 30 ngày:

$A$ = P.t = 10.30.10 = $3000\,\mathrm{kW}$.h

c) Tiền điện của trường học phải trả trong 30 ngày:

Tổng tiền 2000.3000 = 6000000 đồng.

d) Tiền điện của trường học tiết kiệm được trong một năm học:

Tiền điện tiết kiệm = 2000.(10.2.30.9) = 10800000 đồng.
}
\end{ex}

% ===== Câu 6 =====
% ID: L11C4B25-03-TN
% Mức: NB
\begin{ex}
% ID: L11C4B25-03-TN
% Mức: NB
Khi mắc một bóng đèn vào hiệu điện thế $4\,\mathrm{V}$ thl dòng điện qua bóng đèn có cường độ là 600 mA.Công suất tiêu thụ của bóng đèn này là $\nguon{ly11 kntt.pdf}$
\choice
{$24\,\mathrm{W}$.}
{\True $2,4\,\mathrm{W}$.}
{$2400\,\mathrm{W}$.}
{$0,24\,\mathrm{W}$.}
\loigiai{
Chọn B.

Công suất tiêu thụ $P= UI = 4.600\cdot10^{-3} = 2,4\,\mathrm{W}$
}
\end{ex}

% ===== Câu 7 =====
% ID: L11C4B25-04-TL
% Mức: TH
\begin{ex}
% ID: L11C4B25-04-TL
% Mức: TH
Trên một bàn là có ghi $110\,\mathrm{V}$ - $550\,\mathrm{W}$ và trên bóng đèn dây tóc có ghi $110\,\mathrm{V}$ - $100\,\mathrm{W}$. a) Tính điện trở của bàn là và của bóng đèn khi chúng hoạt động bình thường. b) Có thễ mắc nối tiếp bàn là và bóng đèn này vào hiệu điện thế $220\,\mathrm{V}$ được không? Vì sao? (Cho rằng điện trở của bóng đèn và của bàn là không đổi). c) Có thể mắc nối tiếp hai dụng cụ này vào hiệu điện thế lớn nhất là bao nhiêu để chúng không bị hỏng? Tính công suất tiêu thụ của mỗi dụng cụ khi đó. $\nguon{ly11 kntt.pdf}$
\choiceTF

\loigiai{
a) Điện trở của bàn là: $R_{1} = \dfrac{U_{1}^{2}}{\text{P}_{1}} = \dfrac{110^{2}}{550} = 22\Omega$.

Điện trở của bóng đèn: $R_{2} = \dfrac{U_{2}^{2}}{\text{P}_{2}} = \dfrac{110^{2}}{100} = 121\Omega$.

b) Điện trở tương đương của toàn mạch: $R$ = R₁ + R₂ = 22 + 121 = 143 $\Omega$ .

Cường độ dòng điện trong mạch: $I = \dfrac{U}{R} = \dfrac{220}{143} \approx 1,54\text{A}$.

Hiệu điện thế giữa hai đầu bàn là: U₁^(’) = IR₁ = 1, 54.22 = 33, $88\,\mathrm{V}$.

Hiệu điện thế giữa hai đầu bóng đèn: U₂^(’) = I₂ = 1,54.121 = $186\,\mathrm{V}$.

Nhận xét: U₂^(’) > U₂ nên nếu mắc như thế bóng đèn sẽ bị cháy.

c) Cường độ dòng điện định mức của bàn là và của bóng đèn là:

$I_{1} = \dfrac{\text{P}_{1}}{U_{1}} = \dfrac{550}{110} =$ 5 $\mathrm{A}$;I_{2} = \dfrac{\text{P}_{2}}{U_{2}} = \dfrac{100}{110} \approx 0, $91$ \mathrm{A} $.

Khi mắc nối tiếp hai dụng cụ này vào mạch điện, để chúng không bị hỏng thì dòng điện lớn nhất trong mạch có cường độ là I^(') = 0,$91\,\mathrm{A}$.

Hiệu điện thế lớn nhất trong trường hợp này: U^(′) = I^(′)(R₁+R₂) = 0,91.143 = 130,$13\,\mathrm{V}$.

Công suất tiêu thụ trên bàn là: P₁^(′) = I′^{2}.R₁ = 0, 91^{2}.22 ≈ 18,$22\,\mathrm{W}$.

Công suất tiêu thụ trên bóng đèn: P₂^(′) = I′^{2}.R₂ = 0, 91^{2}.121 ≈$ 100\,\mathrm{W}.
}
\end{ex}

% ===== Câu 8 =====
% ID: L11C4B25-04-TN
% Mức: NB
\begin{ex}
% ID: L11C4B25-04-TN
% Mức: NB
Trên một bàn là điện có ghi thông số $220\,\mathrm{V}$ - $1000\,\mathrm{W}$. Điện trở của bàn là điện này là $\nguon{ly11 kntt.pdf}$
\choice
{220 $\Omega$.}
{\True 48,4 $\Omega$.}
{1000 $\Omega$.}
{4,54 $\Omega$.}
\loigiai{
Chọn B.

Điện trở $R = \dfrac{U^{2}}{\text{P}} = \dfrac{220^{2}}{1000} = 48,4\Omega$
}
\end{ex}

% ===== Câu 9 =====
% ID: L11C4B25-05-TL
% Mức: TH
\begin{ex}
% ID: L11C4B25-05-TL
% Mức: TH
Hai dây điện trở của một bếp điện được mắc song song giữa hai điểm $A$ và $B$ có hiệu điện thế $220\,\mathrm{V}$. Cường độ dòng điện qua mỗi dây có giá trị lần lượt là $1,5\,\mathrm{A}$ và 3,5#A. a) Tính điện trở tương đương của đoạn mạch. b) Để có công suất của bếp là $1600\,\mathrm{W}$, người ta phải cắt bỏ bớt một đoạn của dây thứ nhất rồi lại mắc song song với dây thứ hai vào hiệu điện thế nói trên. Hãy tính điện trở của sợi dây bị cắt bỏ đó. $\nguon{ly11 kntt.pdf}$
\choiceTF

\loigiai{
a) Điện trở tương đương của đoạn mạch

$\text{R}_{1} = \dfrac{220}{1,5} \approx 146,7\text{\Omega }\text{R}_{2} = \dfrac{220}{3,5} \approx 63\text{\Omega }\text{R}_{AB} = \dfrac{\text{R}_{1}\text{R}_{2}}{\text{R}_{1} + \text{R}_{2}} = \dfrac{147.63}{147 + 63} \approx 44\text{\Omega }.$

b) Để có công suất là $1600\,\mathrm{W}$ Dây 2 không đổi nên công suất tiêu thụ của dây 2 vẫn là: P₂ = UI₂ = 220.3, 5 = $770\,\mathrm{W}$.

Công suất của dây thứ nhất là: P₁^(′) = $P$ - P₂ = 1600 - 770 = $830\,\mathrm{W}$

Điện trở của dây thứ nhất sau khi cắt: $R_{1}^{'} = \dfrac{U^{2}}{\text{P}_{1}^{'}} = \dfrac{220^{2}}{830} = 58,3\Omega$ Vậy điện trở của sợi dây bị cắt bỏ đó là: R_(cb) = R₁ - R₁^(’) = 146,7 - 58,3 = 88,4 $\Omega$ .
}
\end{ex}

% ===== Câu 10 =====
% ID: L11C4B25-05-TN
% Mức: NB
\begin{ex}
% ID: L11C4B25-05-TN
% Mức: NB
Trên vỏ một máy bơm nước có ghi $220\,\mathrm{V}$ - $1100\,\mathrm{W}$. Cường độ dòng điện định mức của máy bơm là $\nguon{ly11 kntt.pdf}$
\choice
{$I = 0$, $5\,\mathrm{A}$.}
{$I = 50\$, $\mathrm{A}$.}
{\True $I = 5\$, $\mathrm{A}$.}
{$I = 25\$, $\mathrm{A}$.}
\loigiai{
Chọn C.

Cường độ dòng điện $I = \dfrac{\text{P}}{U} = \dfrac{1100}{220} =$ 5 $\mathrm{A}$
}
\end{ex}

% ===== Câu 11 =====
% ID: L11C4B25-06-TL
% Mức: TH
\begin{ex}
% ID: L11C4B25-06-TL
% Mức: TH
Một ấm điện bằng nhôm có khối lượng $0,4 \mathrm{kg}$ chứa $2 \mathrm{kg}$ nước ở 20∘ C.Muốn đun sôi lượng nước đó trong 16 phút thì âm phải có công suất là bao nhiêu? Biết rằng nhiệt dung riêng của nước là $c = 4200\$, $\mathrm{J}/\mathrm{kg}\cdotK$, nhiệt dung riêng của nhôm là $c_1 = 880\$, $\mathrm{J}/\mathrm{kg}\cdotK$ và 27,1% nhiệt lượng toả ra môi trường xung quanh. $\nguon{ly11 kntt.pdf}$
\choiceTF

\loigiai{
Nhiệt lượng cần để tăng nhiệt độ của ấm nhôm từ 20° $C$ tới 100° $C$ là:

Q₁ = m₁c₁(t₂-t₁) = 0,4.880(100-20) = $28160\,\mathrm{J}$

Nhiệt lượng cần để tăng nhiệt độ của nước từ 20° $C$ tới 100° $C$ là:

Q₂ = mc(t₂-t₁) = 2.4200(100-20) = $672000\,\mathrm{J}$

Nhiệt lượng tổng cộng cần thiết: $Q$ = Q₁ + Q₂ = $700160\,\mathrm{J}$ (1)

Mặt khác nhiệt lượng có ích để đun nước do ấm điện cung cấp trong thời gian 16 phút là: $Q$ = HPt (2)

Từ (1) và (2): $\text{P} = \dfrac{Q}{Ht} = \dfrac{700160.100}{(100 - 27,1).960} \approx$ 1000 $\mathrm{W}$.
}
\end{ex}

% ===== Câu 12 =====
% ID: L11C4B25-06-TN
% Mức: NB
\begin{ex}
% ID: L11C4B25-06-TN
% Mức: NB
Nếu đồng thời tăng điện trở dây dẫn, cường độ dòng điện chạy qua dây dẫn lên hai lần, giảm thời gian dòng điện chạy qua dây dẫn hai lần thì nhiệt lượng toả ra trên dây dẫn sẽ tăng $\nguon{ly11 kntt.pdf}$
\choice
{\True 4 lần.}
{8 lần.}
{12 lần.}
{16 lần.}
\loigiai{
Chọn A.

Nhiệt lượng $Q$ = I^{2}Rt , khi đồng thời tăng điện trở dây dẫn, cường độ dòng điện chạy qua dây dẫn lên hai lần, giảm thời gian dòng điện chạy qua dây dẫn hai lần thì nhiệt lượng toả ra trên dây dẫn sẽ là $Q' = \left( {2I} \right)^{2}2R\dfrac{t}{2} = 4I^{2}Rt$ nên tăng 4 lần.
}
\end{ex}

% ===== Câu 13 =====
% ID: L11C4B25-07-TL
% Mức: TH
\begin{bt}
% ID: L11C4B25-07-TL
% Mức: TH
Một bếp điện sợi đốt tiêu thụ công suất $P = 1$, $1 \mathrm{kW}$ được dùng ở mạng điện có hiệu điện thế $U = 120\$, $\mathrm{V}$. Dây nối từ ổ cắm vào bếp điện $r = 1\Omega$. a) Tính điện trở $R$ của bếp điện khi hoạt động bình thường. b) Tính nhiệt lượng toả ra ở bếp điện khi sử dụng liên tục bếp điện trong thời gian nửa giờ. $\nguon{ly11 kntt.pdf}$
\loigiai{
a) Áp dụng công thức tính công suất tiêu thụ:

$P= I^{2}R$ với $\left. I = \dfrac{U}{R + r}\Rightarrow\text{P} = \dfrac{U^{2}}{(R + r)^{2}} \cdot R \right.\Rightarrow$ 120^{2}R=1100 $(R^{2}+2R+1)$ ⇔11R^{2}-122R+11=0

$\left. \Rightarrow R = 11\text{\Omega };R = \dfrac{1}{11}\text{\Omega }. \right.$

Loại nghiệm $R$ = $\dfrac{1}{11}\Omega$ vì khi đó hiệu điện thế trên điện trở bằng $10\,\mathrm{V}$, hiệu điện thế trên điện trở dây bằng $110\,\mathrm{V}$, dẫn tới công suất toả nhiệt trên dây nối quá lớn, không thực tế.

b) Nhiệt lượng toả ra trên bếp điện trong thời gian nửa giờ $Q$ = Pt = 1980kJ.
}
\end{bt}

% ===== Câu 14 =====
% ID: L11C4B25-07-TN
% Mức: NB
\begin{ex}
% ID: L11C4B25-07-TN
% Mức: NB
Công suất điện cho biết $\nguon{ly11 kntt.pdf}$
\choice
{khả năng thực hiện công của dòng điện.}
{năng lượng của dòng điện.}
{\True lượng điện năng sử dụng trong một đơn vị thời gian.}
{mức độ mạnh - yếu của dòng điện.}
\loigiai{
Chọn C.

Công suất điện cho biết lượng điện năng sử dụng trong một đơn vị thời gian.
}
\end{ex}

% ===== Câu 15 =====
% ID: L11C4B25-08-TL
% Mức: TH
\begin{bt}
% ID: L11C4B25-08-TL
% Mức: TH
Nguồn điện có điện trở trong $r = 2\Omega$, cung cấp một công suất $P$ cho mạch ngoài là điện trở $R_1 = 0$,5 $\Omega$. Mắc thêm vào mạch ngoài điện trở R_2 thì công suất tiêu thụ mạch ngoài không đổi. Hỏi R_2 nối tiếp hay song song với R_1 và có giá trị bao nhiêu? $\nguon{ly11 kntt.pdf}$
\loigiai{
Khi mạch ngoài chỉ có điện trở R₁ thì công suất tiêu thụ mạch ngoài:

$\text{P} = \left( \dfrac{\text{E}}{R_{1} + r} \right)^{2}.R_{1}$ (1)

Nếu mắc thêm điện trở R₂ thì điện trở mạch ngoài là R₁₂.

Theo đầu bài, ta có $\text{P} = \dfrac{\text{E}^{2}}{\left( {R_{12} + r} \right)^{2}} \cdot R_{12}$ (2)

Từ (1) và (2): 0,5(R₁₂+2)^{2}=R₁₂(0,5+2)^{2} $\Rightarrow$ R₁₂^{2}-8,5R₁₂+4=0.

Giải phương trình ta thu được: R₁₂ = 8 $\Omega$ .

R₁₂ = 0, 5 $\Omega$  (loại vì R₁₂ ≠ R₁)

Nhận thấy R₁₂ > R₁ nên R₂ phải mắc nối tiếp với R₁ và R₂ = R₁₂ - R₁ = 7, 5 $\Omega$ .
}
\end{bt}

% ===== Câu 16 =====
% ID: L11C4B25-08-TN
% Mức: NB
\begin{ex}
% ID: L11C4B25-08-TN
% Mức: NB
Trên các thiết bị điện gia dụng thường có ghi $220\,\mathrm{V}$ và số oát (W). Số oát này có ý nghĩa gì? $\nguon{ly11 kntt.pdf}$
\choice
{Công suất tiêu thụ điện của dụng cụ khi nó được sử dụng với những hiệu điện thế nhỏ hơn $220\,\mathrm{V}$.}
{\True Công suất tiêu thụ điện của dụng cụ khi nó được sử dụng với đúng hiệu điện thế $220\,\mathrm{V}$.}
{Công mà dòng điện thực hiện trong một phút khi dụng cụ này được sử dụng với đúng hiệu điện thế $220\,\mathrm{V}$.}
{Điện năng mà dụng cụ tiêu thụ trong một giờ khi nó được sử dụng với đúng hiệu điện thế $220\,\mathrm{V}$}
\loigiai{
Chọn B.

Số oát có ý nghĩa là công suất tiêu thụ điện của dụng cụ khi nó được sử dụng với đúng hiệu điện thế $220\,\mathrm{V}$.
}
\end{ex}

% ===== Câu 17 =====
% ID: L11C4B25-09-TL
% Mức: TH
\begin{ex}
% ID: L11C4B25-09-TL
% Mức: TH
Một nguồn điện có suất điện động $\mathcal{E} = 6\$, $\mathrm{V},$ điện trở trong $r = 2\Omega$, mạch ngoài có điện trở R. a) Tính $R$ để công suất tiêu thụ ở mạch ngoài là $4\,\mathrm{W}$. b) Với giá trị nào của $R$ thì công suất tiêu thụ ở mạch ngoài là lớn nhất. Tính giá trị đó. $\nguon{ly11 kntt.pdf}$
\choiceTF

\loigiai{
a) Công suất tiêu thụ ở mạch ngoài $\text{P}_{n} = \left( \dfrac{\text{E}}{r + R} \right)^{2}R\left. \Rightarrow 4 = \left( \dfrac{6}{2 + R} \right)^{2} \cdot R\Rightarrow 4\left( {4 + 4R + R^{2}} \right) = 36R \right.\Rightarrow$ R^{2}-5R+4=0. Giải phương trình thu được: $R$ = 1 $\Omega$ ; $R$ = 4 $\Omega$ ..

b) Biến đổi, đưa công thức tính công suất P_(n) về dạng $\text{P}_{n} = \left( \dfrac{\text{E}}{\dfrac{r}{\sqrt{R}} + \sqrt{R}} \right)^{2}$ Để P_(n)max thì $\left( {\sqrt{\text{R}} + \dfrac{\text{r}}{\sqrt{\text{R}}}} \right)\text{min}$

Theo bất đẳng thức Cauchy, ta có: $\left. \sqrt{R} = \dfrac{r}{\sqrt{R}}\Rightarrow R = r = 2\Omega \right.$

Khi đó $\text{P}_{n} = \dfrac{36}{8} = 4,$ 5 $\mathrm{W}$.
}
\end{ex}

% ===== Câu 18 =====
% ID: L11C4B25-09-TN
% Mức: NB
\begin{ex}
% ID: L11C4B25-09-TN
% Mức: NB
Công suất định mức của các dụng cụ điện là $\nguon{ly11 kntt.pdf}$
\choice
{công suất lớn nhất mà dụng cụ đó có thể đạt được.}
{công suất tối thiểu mà dụng cụ đó có thể đạt được.}
{\True công suất đạt được khi nó hoạt động bình thường.}
{công suất trung bình của dụng cụ đó.}
\loigiai{
Chọn C.

Công suất định mức của các dụng cụ điện là công suất đạt được khi nó hoạt động bình thường.
}
\end{ex}

% ===== Câu 19 =====
% ID: L11C4B25-10-TL
% Mức: TH
\begin{ex}
% ID: L11C4B25-10-TL
% Mức: TH
Một dây điện trở có thể làm 500 mL nước tăng nhiệt độ thêm 60∘ trong 5 phút khi hoạt động ở hiệu điện thế $220\,\mathrm{V}$. Biết nhiệt dung riêng của nước là $4180\,\mathrm{J}$ / $\mathrm{kg}$. $K$ a) Tính công suất toả nhiệt của điện trở. b) Cường độ dòng điện qua điện trở là bao nhiêu? $\nguon{ly11 kntt.pdf}$
\choiceTF

\loigiai{
Nhiệt lượng nước thu được: $Q$ = mc.$\Delta$ t - 0,5.4.18$\cdot$ 10^{3}.60 = $125400\,\mathrm{J}$.

a) Nhiệt lượng trên chính là điện năng chuyển hoá thành, như vậy điện năng tiêu thụ trên điện trở có giá trị $A$ = $125400\,\mathrm{J}$.

Do đó ta suy ra công suất tiêu thụ trên điện trở: $\text{P} = \dfrac{A}{t} = \dfrac{125400}{5.60} =$418 $\mathrm{W}$.

b) Cường độ dòng điện: $I$ = $\dfrac{\text{P}}{U}$ = $\dfrac{418}{220}$ = 1,$9\,\mathrm{A}$.
}
\end{ex}

% ===== Câu 20 =====
% ID: L11C4B25-10-TN
% Mức: TH
\begin{ex}
% ID: L11C4B25-10-TN
% Mức: TH
Công thức nào trong các công thức sau đây cho phép xác định năng lượng điện tiêu thụ của đoạn mạch (trong trường hợp dòng điện không đổi)? UI $\nguon{ly11 kntt.pdf}$
\choice
{$A = UI 2 t$.}
{$A = U 2 It$.}
{\True $A = Ult$.}
{$A = t$.}
\loigiai{
Chọn C.

Công thức đúng $A$ = UIt.
}
\end{ex}

% ===== Câu 21 =====
% ID: L11C4B25-11-TL
% Mức: TH
\begin{ex}
% ID: L11C4B25-11-TL
% Mức: TH
Hai điện trở $R_1 = 20\Omega và điện trở R_2 chưa biết giá trị$ được mắc nối tiếp với nhau và mắc vào hiệu điện điện thế $U = 220\$, $\mathrm{V}$ thì điện trở R_2 tiêu thụ một công suất là $P_2 = 600\$, $\mathrm{W}$. Tính giá trị điện trở R_2 biết rằng dòng điện chạy qua điện trở R_2 có giá trị không lớn hơn $5\,\mathrm{A}$. $\nguon{ly11 kntt.pdf}$
\choiceTF

\loigiai{
Gọi $I$ là dòng điện chạy qua đoạn mạch, ta có: $P$ = P₁ + P₂

UI = R₁I^{2} + R₂I^{2}

Thay các số liệu đầu bài cho, ta có: R₁I^{2}-UI+600=0⇔20I^{2}-220I+600=0

[Hai điện trở R1 = 20 $\Omega$  và điện trở R2 chưa biết giá trị]. Vì R₁ mắc nối tiếp với R₂ nên: I₁ = I₂ = I.

Theo đề bài $. I_{2} \leq$ 5 $\mathrm{A}$ \Rightarrow I_{2} = $5$ \mathrm{A} $\Rightarrow R_{2} = \dfrac{\text{P}}{I_{2}^{2}} = \dfrac{600}{5^{2}} = 24\Omega .$
}
\end{ex}

% ===== Câu 22 =====
% ID: L11C4B25-11-TN
% Mức: TH
\begin{ex}
% ID: L11C4B25-11-TN
% Mức: TH
Công thức nào dưới đây không phải là công thức tính công suất của vật tiêu thụ điện toả nhiệt? U_2 $\nguon{ly11 kntt.pdf}$
\choice
{$P = UI$.}
{$P = RI_2$.}
{\True $P = R_2$.}
{$P = R$.}
\loigiai{
Chọn C.

Công thức sai $P$ = IR^{2}
}
\end{ex}

% ===== Câu 23 =====
% ID: L11C4B25-12-TL
% Mức: TH
\begin{ex}
% ID: L11C4B25-12-TL
% Mức: TH
Từ một nguồn điện có hiệu điện thế $U$, điện năng được truyền qua dây dẫn tới nơi tiêu thụ. Biết điện trở dây dẫn là $R = 5\Omega$, công suất của nguồn điện là $P =$ 62 $\mathrm{kW}$. Tìm độ giảm thế trên dây, công suất hao phí và hiệu suất tải điện nếu: a) $U = 6200\$, $\mathrm{V}$. b) $U = 620\$, $\mathrm{V}$. $\nguon{ly11 kntt.pdf}$
\choiceTF

\loigiai{
a) $\Delta U =$ 50 $\mathrm{V}$;\Delta\text{P} = $500$ \mathrm{W} $;H = \dfrac{\text{P} - \Delta\text{P}}{\text{P}} \approx 99,2\%$.

b) $\Delta U =$ 500 $\mathrm{V}$;\Delta\text{P} = $50000$ \mathrm{W} $;H = \dfrac{\text{P} - \Delta\text{P}}{\text{P}} \approx 19,35\%$.
}
\end{ex}

% ===== Câu 24 =====
% ID: L11C4B25-13-TL
% Mức: TH
\begin{bt}
% ID: L11C4B25-13-TL
% Mức: TH
Người ta dùng một ấm nhôm có khối lượng $m_1 = 0$, $4 \mathrm{kg}$ để đun một lượng nước $m_2 =$ 2 $\mathrm{kg}$ thì sau 20 phút nước sẽ sôi $. Bếp điện có hiệu suất$ H = 60% và $được dùng ở mạng điện có hiệu điện thế$ U = 220\ $,$ \mathrm{V} $. Nhiệt độ ban đầu của nước là$ t_1 = 20∘ C $, nhiệt dung riêng của nhôm là$ c_1 = 920\ $,$ \mathrm{J}/\mathrm{kg} $.K$, của nước là $c_2 = 4200\$, $\mathrm{J}/\mathrm{kg}$.K. Tìm nhiệt lượng cần cung cấp cho âm nước và cường độ dòng điện chạy qua bếp điện. $\nguon{ly11 kntt.pdf}$
\loigiai{
Nhiệt lượng cần cung cấp cho ấm nước để đun sôi lượng nước đã cho là:

$Q= (c₁m₁+c₂m₂)(t₂-t₁) =701440\,\mathrm{J}$.

Do hiệu suất của bếp nhỏ hơn 100% nên nhiệt lượng do bếp toả ra là:

$Q_{\text{bep}} = \dfrac{Q}{H} = 1169066,7\text{J}$

Mặt khác $\left. Q_{\text{b}e\text{p}} = UIt\Rightarrow I = \dfrac{Q_{\text{bep}}}{Ut} = 4,4\text{A} \right.$
}
\end{bt}
